% \documentclass[10pt,respuestas,a4]{examenEPN}
\documentclass[10pt,a4]{examenEPN0}
\usepackage[T1]{fontenc}
\examen{Prueba 01}
\usepackage{comandosEPN}

\usepackage{empheq}

\newcommand{\itemcp}{\stepcounter{enumi}\item[ \textcolor{red}{\theenumi.$^*$}] }

\newcommand{\iitemcp}{\stepcounter{enumii}\item[ \textcolor{red}{\theenumii)$^*$}] }


\materia{Álgebra Lineal}
%\autor{Departamento de Formación Básica}
\autor{GR \hspace{1cm}.}
%\fecha{Semestre 2026-A}
\fecha{Nombre: \hspace{9cm} Cédula:}
\logouno[1.6cm]{Logo_EPN.eps}
\logodos[2.6cm]{Logo_DFB.eps}

\newcommand{\la}{\lambda}

%Datos
\newcommand{\datosAlgebra}[2]{
\EnExamen
{
\vspace{5mm}
    \noindent
    Nombres: \hrulefill\
    Carrera: \rule{4.5cm}{0.4pt}\
    \ifthenelse{\equal{#1}{}}%
    {Paralelo: \underline{\parbox{1.8cm}{\parbox{1cm}{}}}\\[3mm]}
    {Paralelo: \underline{\parbox{1.8cm}{\centering #1}}\\[3mm]}
    C\'edula: \rule{5cm}{0.4pt}\
    Firma: \hrulefill\
    Fecha: \underline{\parbox{2.1cm}{\centering #2}}
}}


% \newcommand{\indicacionesPrueba}{%
% \EnExamen{
% \begin{indicaciones}
%     \item La prueba tiene una duraci\'on  de \textbf{120 minutos.}
%     \item No se permite la utilizaci\'on de dispositivos electr\'onicos (celular, tablets, calculadora, etc.)
%     \item Leer detenidamente los enunciados, \textbf{justificar la respuesta para acceder al puntaje total}.
%     \item La prueba es completamente individual. Los estudiantes deberán dejar sus teléfonos celulares guardados y no podrán tenerlos consigo (por ejemplo, en los bolsillos de su ropa) durante toda la evaluación. En caso de detectarse algún acto de deshonestidad académica o el incumplimiento de esta norma, no podrán seguir rindiendo la prueba y su calificación será \textbf{cero}.
% \end{indicaciones}
% \pagestyle{empty}
% }}

\newcommand{\indicacionesPrueba}{%
\EnExamen{
\begin{indicaciones}
    \item La prueba tiene una duraci\'on  de \textbf{60 minutos.}
    \item Leer detenidamente los enunciados, \textbf{justificar la respuesta} para acceder al puntaje total.
    \item No se permite la utilizaci\'on de dispositivos electr\'onicos (celular, tablets, calculadora, etc.)
    \item La prueba es completamente individual
    \item No se permite portar teléfonos celulares, relojes inteligentes ni otros dispositivos electrónicos durante el examen, aun cuando permanezcan en los bolsillos. Estos dispositivos deberán guardarse en la mochila antes del inicio de la evaluación y esta deberá mantenerse cerrada durante todo el examen.
    \item He leído las indicaciones y acepto que el incumplimiento de estas se penaliza con \textbf{cero}. Firma:...................................

    
\end{indicaciones}
\pagestyle{empty}
}}

\newcommand{\indicacionesExamen}{
\EnExamen{\begin{indicaciones}
    \item La evaluación tiene una duraci\'on  de \textbf{120 minutos.}
    \item No se permite la utilizaci\'on de dispositivos electr\'onicos (celular, tablets, calculadora, etc.)
    %\item Leer detenidamente los enunciados, justificar la respuesta para acceder al puntaje total.
    \item El examen es completamente individual. En caso de detectarse alg\'un acto de deshonestidad acad\'emica no podr\'a seguir rindiendo la prueba y su calificaci\'on ser\'a \textbf{cero}.
\end{indicaciones}
}\pagestyle{empty}
}

%% Imprimir las una hoja A4, dos veces en A5
\newcommand{\imprimirDoble}
{%
\usepackage{pgfmorepages}
\pgfmorepagesloadextralayouts
\pgfpagesuselayout{repeated 2-up}[a4paper,landscape]
\thispagestyle{empty}
}


%Este sí se puede comentar
\newenvironment{solucion}{\color{blue}\proof[Solución]}{\endproof\color{black}}




% \renewcommand{\@autor}{
% ---\\
% \phantom{.} 
% ...}
% \makeatother

\begin{document}
\encabezado
\datosAlgebra{GR6M}{23/10/2026}
%Poner el grupo y la fecha, o dejar las llaves vacías.
\indicacionesPrueba
\begin{preguntas}
\item Responda:
\begin{enumerate}
    \item Sean \(A\in \mathbb R^{3\times 4}, B\in \mathbb R^{3\times 2}\). Determine la cantidad de filas y columnas que deben tener las matrices \(C\) y \(D\) para que la siguiente ecuación tenga sentido: 
    \(A^TB = DC^TB+D\) \puntaje{2}
    \item Calcule \(\text{Tr}(B^TA)\) donde:\puntaje{3}
    \[A = \left(\begin{array}{cc}
         -3&5  \\
         4&0 
    \end{array}\right),\quad B = \left(\begin{array}{cc}
         1&0  \\
         2&-1 
    \end{array} \right).\]
\end{enumerate}
% Espacio de respuesta hasta el final de la página.
\par\vfill\newpage

%%%%%%%%%%%%%%%%%%%%%%%%%%%%%%%%%%


\item Plantee un sistema lineal de la forma \(Ax = b\) que permita modelar la siguiente situación (No es necesario resolver el sistema, pero sí se debe especificar la matriz \(A\) y el vector \(b\)).

\textit{Se llevó a cabo un concierto en un estadio con capacidad para \(21.000\) personas, donde todas las entradas fueron vendidas. Según la localidad, los boletos tenían un costo de: General (\(50\$\)), Tribuna (\(75\$\)), VIP (\(150\$\)), y la recaudación total fue de \(1.375.000\$\) . La cantidad de asistentes a la localidad general es el doble de asistentes de las otras localidades juntas. ¿Cuántas personas asistieron a cada localidad?
}
\puntaje{5}
% Espacio de respuesta hasta el final de la página.
\par\vfill\newpage



%%%%%%%%%%%%%%%%%%%%%%%%%%%%%%%%%

\item Considere el siguiente sistema de ecuaciones lineales, donde \(\alpha\) es un parámetro:
\[
\begin{aligned} 
x - (2\alpha)y + 4z &= 6 \\ 
-x + (12+4\alpha)y  -12z &= -10 \\ 
 -3x-36y 
 + (\alpha^2+\alpha-18)z &= 6\alpha+30
\end{aligned}
\]

\begin{enumerate}
    \item Halle la matriz ampliada asociada al sistema y utilice los siguientes pasos para escalonarla.  \puntaje{1}
    \begin{itemize}
        \item \( F_1+F_2 \rightarrow F_2 \)
        \item \( 3F_1+F_3 \rightarrow F_3 \)
        \item \( 3F_2+F_3 \rightarrow F_3 \)
    \end{itemize}

    \item ¿Para qué valor de \(\alpha\) el sistema no tiene solución y para qué valor tiene infinitas soluciones? En el caso de infinitas soluciones, halle el conjunto solución. \puntaje{4}
    
\end{enumerate}
% Espacio de respuesta hasta el final de la página.
\par\vfill\newpage

%%%%%%%%%%%%%%%%%%%%%%%%%%%%%%%%%%%%%%%

\item (Ejercicio del deber) Dado el siguiente sistema:
\[
\left\{
\begin{aligned}
x + y + z &= 6 \\
2x - y + 3z &= 9 \\
x + 2y - z &= 2
\end{aligned}
\right.
\]

 Resuelva el sistema utilizando el método de \textbf{Gauss}. Determine $\operatorname{rang}(A)$ y $\operatorname{rang}(A|b)$. \puntaje{5}



%%%%%%%%%%%%%%%%%%%%%%%%%%%%%%%%%%%%%%%%%

    
% El espacio restante de la última página queda disponible para responder.
\par\vfill
\end{preguntas}
\final
\end{document}
